%%%%%%%%%%%%%%%%%%%%%%%%%%%%%%%%%%%%%%%%%%%%%%%%%%%%%%%%%%%%%%%%%%%%%%
% Lớp 11 | Hàm số lượng giác và phương trình lượng giác | [Giảm] Phương trình lượng giác thường gặp
% Dạng 1D1?6-3: Phương trình bậc nhất đối với sinx và cosx
% Số câu: 2  (sinh tự động bởi tach_cau.py; từ nay thầy sửa trực tiếp file này)
%%%%%%%%%%%%%%%%%%%%%%%%%%%%%%%%%%%%%%%%%%%%%%%%%%%%%%%%%%%%%%%%%%%%%%

%%%%%%%%%%%%%%%% ĐÚNG/SAI %%%%%%%%%%%%%%%%

% nguồn: Ok_dungsai_lop11_Lan1
\begin{ex}%[1D1H6-3]%[ma:bf62b1b9]
 Xét tính đúng sai của các mệnh đề sau. 
 \choiceTF 
  {Phương trình $a\sin x+b\cos x=c$ có nghiệm khi  $a^2+b^2<c^2$}
  {\True Phương  trình $a\sin x+b\cos x=c$ có nghiệm khi $a^2+b^2\geq c^2+1$}
  {\True Phương trình $a\sin x+b\cos x=c$ có nghiệm khi $a^2+b^2\geq c^2$} 
   {\True Phương trình $a\sin x+b\cos x=c$ có nghiệm khi $a^2+b^2>c^2$}
 \loigiai{Phương trình  $a\sin x+b\cos x=c$ có nghiệm khi và chỉ khi $a^2+b^2\geq c^2$.\\
 Ta có $a^2+b^2\ge c^2+1\Rightarrow a^2+b^2\ge c^2$, suy ra phương trình $a\sin x+b\cos x=c$  có nghiệm khi $a^2+b^2\geq c^2+1$.}
\end{ex}

% nguồn: Ok_dungsai_lop11_Lan1
\begin{ex}%[1D1H6-3]%[ma:f21cade2]
 Phương trình  $\sqrt{3}\sin x+\cos x=1$ tương đương với phương trình nào sau đây?
  \choiceTF
   {$\sin\left(x+\dfrac{\pi}{3}\right)=\dfrac{1}{2}$}
  {\True $\cos\left(x-\dfrac{\pi}{3}\right)=\dfrac{1}{2}$} 
  {\True  $\sin\left(x+\dfrac{\pi}{6}\right)=\dfrac{1}{2}$}
  {$\cos\left(x-\dfrac{\pi}{6}\right)=\dfrac{1}{2}$} 
 \loigiai{ 
  Chia hai vế của  phương trình cho $2$, ta được
  \begin{eqnarray*}
   &\sqrt{3}\sin x+\cos x=1&\Leftrightarrow\dfrac{\sqrt{3}}{2}\sin  x+\dfrac{1}{2}\cos x=\dfrac{1}{2}\\
   &&\Leftrightarrow\sin\dfrac{\pi}{3}\sin  x+\cos\dfrac{\pi}{3}\cos x=\dfrac{1}{2}\\
   &&\Leftrightarrow\cos\left(x-\dfrac{\pi}{3}\right)=\dfrac{1}{2}\\ 
   &&\Leftrightarrow\sin\left(x+\dfrac{\pi}{6}\right)=\dfrac{1}{2}. 
  \end{eqnarray*}
  }
\end{ex}
